\subsection{Uniform Integrability and Vitali Convergence}

Uniform integrability is the $L^1$ condition that rules out concentration of mass on tiny sets. On infinite-measure spaces it must be paired with tightness, which rules out mass escaping to remote finite-measure tails.

\begin{lemma} \label{lem:finite-measure-lebesgue-small-piece-decomposition} \hlt{Finite-Measure Lebesgue Sets Have Small Pieces}\\
Let $E\in\mathcal{L}^1$, $m(E)<+\infty$. $\forall \delta>0$, $\exists$ pairwise disjoint $E_1,\ldots,E_N\in\mathcal{L}^1$ such that $E=\bigsqcup\limits_{k=1}^NE_k$, $m(E_k)<\delta$ $\forall k$.
\begin{proof} Choose $n_0$ with $T\defeq E\setminus[-n_0,n_0]$ satisfying $m(T)<\delta$.\\
Partition $[-n_0,n_0]$ into finitely many disjoint intervals $I_j$ with $m(I_j)<\delta$.\\
Discarding empty sets, the pieces $T$ and $E\cap I_j$ give the required finite disjoint decomposition.\\
\hlt{Structural Commentary}: Finite measure gives a small tail $T$; a fine partition makes $E\setminus T$ small-piece.\\
\hlt{Key Insight}: One small tail plus a fine finite grid gives a finite small-piece decomposition.
\end{proof}
\end{lemma}

\begin{proposition} \label{prop:l1-small-set-absolute-continuity-finite-lebesgue-converse} \hlt{$L^1$ Absolute Continuity on Small Sets}\\
Let $E\in\mathcal{L}^1$, let $\mathcal{L}^1_E=\{A\in\mathcal{L}^1:A\subseteq E\}$, let $m_E=m|_{\mathcal{L}^1_E}$, and let $f:E\to\R$ be measurable.\\
Consider the condition $(\ast)$: $\forall \varepsilon>0$ $\exists \delta>0$ such that $A\in\mathcal{L}^1_E$ and $m(A)<\delta$ imply $\int_A\abs f\,dm<\varepsilon$.
\begin{enumerate}[label=(\roman*)]
\setlength{\itemsep}{0pt}
\item If $f\in L^1(m_E)$, then $(\ast)$ holds.
\item If $m(E)<+\infty$ and $(\ast)$ holds, then $f\in L^1(m_E)$.
\end{enumerate}
\begin{proof}
\begin{enumerate}[label=(\roman*)]
\setlength{\itemsep}{0pt}
\item Lemma \ref{lem:lp-integral-absolute-continuity} on $(E,\mathcal{L}^1_E,m_E)$ with $p=1$ gives $f\in L^1(m_E)\Rightarrow(\ast)$.
\item Choose $\delta>0$ from $(\ast)$ with $\varepsilon=1$. By Lemma \ref{lem:finite-measure-lebesgue-small-piece-decomposition}, write $E=\bigsqcup_{k=1}^{N}E_k$ with $m(E_k)<\delta$.\\
Then $\int_E\abs f\,dm=\sum_{k=1}^{N}\int_{E_k}\abs f\,dm<N$, so $f\in L^1(m_E)$.
\end{enumerate}
\hlt{Structural Commentary}: $L^1$ gives small-set control; finite partitions recover $L^1$.\\
\hlt{Key Insight}: A finite small-piece decomposition turns local integral bounds into one global finite integral.
\end{proof}
\end{proposition}

\begin{definition} \label{def:l1-uniform-integrability-and-tightness} \hlt{Uniform Integrability and Tightness}\\
Let $\mathcal{F}$ be a family of measurable functions $h:X\to\R$ on $(X,\mathcal{M},\mu)$.
\begin{enumerate}[label=(\roman*)]
\setlength{\itemsep}{0pt}
\item $\mathcal{F}$ is \hlt{uniformly integrable} on $X$ if $\forall\varepsilon>0\,\exists\delta>0$ such that $A\in\mathcal{M}$ and $\mu(A)<\delta$ imply $\sup_{h\in\mathcal{F}}\int_A\abs h\,d\mu<\varepsilon$.
\item $\mathcal{F}$ is \hlt{tight} on $X$ if $\forall\varepsilon>0\,\exists X_\varepsilon\in\mathcal{M}$ with $\mu(X_\varepsilon)<+\infty$ and $\sup_{h\in\mathcal{F}}\int_{X\setminus X_\varepsilon}\abs h\,d\mu<\varepsilon$.
\end{enumerate}
On finite measure spaces, tightness is automatic by taking $X_\varepsilon=X$.\\
\hlt{Key Insight}: UI controls mass on small $A$; tightness controls mass on $X\setminus X_\varepsilon$.
\end{definition}

\begin{example} \label{ex:uniform-integrability-tightness-obstructions} \hlt{Concentration versus Escape}\\
Figure \ref{fig:fatou-concentration-escape} separates the two obstructions.
\begin{enumerate}[label=(\roman*)]
\setlength{\itemsep}{0pt}
\item \hlt{Concentration}: on $(0,1)$, $u_n\defeq n\mathbf{1}_{(0,1/n)}$ is not uniformly integrable.\\
Given $\delta>0$, choose $n>1/\delta$ and $A=(0,1/n)$, so $m(A)<\delta$ but $\int_Au_n\,dm=1$.
\item \hlt{Escape}: on $\R$, $v_n\defeq\mathbf{1}_{I_n}$ with $I_n=(n,n+1)$ is uniformly integrable since $\int_Av_n\,dm\leq m(A)$, but not tight: if $m(E_0)<+\infty$ and $\int_{\R\setminus E_0}v_n\,dm<1/2$ for all $n$, then $m(E_0\cap I_n)>1/2$ for all $n$, contradiction.
\end{enumerate}
\hlt{Structural Commentary}: Concentration occurs even on $m(X)<+\infty$; escape is an $m(X)=+\infty$ obstruction.\\
\hlt{Key Insight}: Uniform integrability controls small-set mass; tightness keeps mass in finite-measure cores.
\end{example}

\begin{proposition} \label{prop:l1-uniform-integrability-basic-sources} \hlt{Basic Sources of Uniform Integrability and Tightness}\\
Let $(X,\mathcal{M},\mu)$ be a measure space.
\begin{enumerate}[label=(\roman*)]
\setlength{\itemsep}{0pt}
\item If $0\leq g\in L^1(\mu)$ and $\abs h\leq g$ a.e. for all $h\in\mathcal{F}$, then $\mathcal{F}$ is uniformly integrable and tight.
\item If $f_1,\ldots,f_N\in L^1(\mu)$, then $\{f_1,\ldots,f_N\}$ is uniformly integrable and tight.
\end{enumerate}
\begin{proof} Let $\varepsilon>0$.
\begin{enumerate}[label=(\roman*)]
\setlength{\itemsep}{0pt}
\item Lemma \ref{lem:lp-integral-absolute-continuity} gives $\delta>0$ with $\mu(A)<\delta\Rightarrow\int_Ag\,d\mu<\varepsilon$, hence $\sup_{h\in\mathcal F}\int_A\abs h\,d\mu<\varepsilon$.\\
Proposition \ref{prop:finite-positive-integral-localization} gives finite-measure $X_\varepsilon$ with $\int_{X\setminus X_\varepsilon}g\,d\mu<\varepsilon$, hence $\sup_{h\in\mathcal F}\int_{X\setminus X_\varepsilon}\abs h\,d\mu<\varepsilon$.
\item Apply (i) to each $\{f_j\}$. Take $\delta=\min\limits_j\delta_j$ for small-set bounds, $X_\varepsilon=\bigcup\limits_{j=1}^NX_{\varepsilon,j}$ for tail bounds.
\end{enumerate}
\hlt{Structural Commentary}: An envelope gives common bounds; finite families use $\min_j\delta_j$ and $\bigcup_jX_{\varepsilon,j}$.\\
\hlt{Key Insight}: Domination gives uniform estimates; finite families allow minima and finite unions.
\end{proof}
\end{proposition}

\begin{theorem} \label{thm:vitali-convergence-l1-ui-tight-ae} \hlt{Vitali Convergence for $L^1$ Families}\\
Let $(X,\mathcal{M},\mu)$ be a measure space, let $f_n\in L^1(\mu)$, let $f:X\to\R$ be measurable, and assume $f_n\to f$ a.e.
\begin{enumerate}[label=(\roman*)]
\setlength{\itemsep}{0pt}
\item If $\mu(X)<+\infty$ and $\{f_n:n\geq1\}$ is uniformly integrable, then $f\in L^1(\mu)$ and $\norm{f_n-f}_1\to0$.
\item If $\{f_n:n\geq1\}$ is uniformly integrable and tight, then $f\in L^1(\mu)$ and $\norm{f_n-f}_1\to0$.
\end{enumerate}
Consequently, $\int_Xf_n\,d\mu\to\int_Xf\,d\mu$.
\begin{proof}
\begin{enumerate}[label=(\roman*)]
\setlength{\itemsep}{0pt}
\item Corollary \ref{cor:finite-measure-ae-implies-measure} gives $f_n\to f$ in measure. Apply Theorem \ref{thm:vitali-convergence-lp} with $p=1$: UI is the small-set condition, and the finite-tail condition holds with $E_\varepsilon=X$.
\item First prove $f_n\to f$ in measure. Fix $\alpha,\eta>0$ and set $B_n^\alpha=\{\abs{f_n-f}\geq\alpha\}$.\\
By tightness, choose $X_0\in\mathcal{M}$ with $\mu(X_0)<+\infty$ and $\sup_n\int_{X\setminus X_0}\abs{f_n}\,d\mu<\alpha\eta/4$; Fatou gives the same bound for $f$. Hence $\mu(B_n^\alpha\cap(X\setminus X_0))<\eta/2$. On $X_0$, Corollary \ref{cor:finite-measure-ae-implies-measure} gives $\mu(B_n^\alpha\cap X_0)<\eta/2$ for all large $n$. Thus $f_n\to f$ in measure on $X$. Apply Theorem \ref{thm:vitali-convergence-lp} with $p=1$, using UI and tightness.
\end{enumerate}
The integral convergence follows from $\abs{\int_X(f_n-f)\,d\mu}\leq\norm{f_n-f}_1$.\\
\hlt{Structural Commentary}: Tightness gives finite cores; UI upgrades in-measure convergence to $L^1$.\\
\hlt{Key Insight}: Uniform integrability prevents concentration; tightness prevents escape.
\end{proof}
\end{theorem}

\begin{corollary} \label{cor:nonnegative-l1-vanishing-vitali-criterion} \hlt{Nonnegative Integral Vanishing Criterion}\\
Let $f_n\in L^1(\mu)$ with $f_n\geq0$ a.e. Then $\int_X f_n\,d\mu\to0$ iff:
\begin{enumerate}[label=(\roman*)]
\setlength{\itemsep}{0pt}
\item $f_n\to0$ in measure;
\item $\{f_n:n\geq1\}$ is uniformly integrable;
\item $\{f_n:n\geq1\}$ is tight.
\end{enumerate}
\begin{proof}
\begin{enumerate}[label=\roman*.]
\setlength{\itemsep}{0pt}
\item[$(\Rightarrow)$] Since $f_n\geq0$, $\norm{f_n}_1=\int_Xf_n\,d\mu\to0$. Thus $f_n\to0$ in measure, and $\{f_n:n\geq1\}$ is uniformly integrable and tight by Theorem \ref{thm:vitali-convergence-lp} with $p=1$.
\item[$(\Leftarrow)$] Conditions (i)--(iii) are exactly the in-measure, small-set, and finite-tail hypotheses of Theorem \ref{thm:vitali-convergence-lp} with $p=1$ and limit $0$. Hence $\norm{f_n}_1\to0$, so $\int_Xf_n\,d\mu\to0$.
\end{enumerate}
\hlt{Structural Commentary}: For nonnegative functions, integral convergence to zero is $L^1$ convergence to zero.\\
\hlt{Key Insight}: Vitali's conditions characterize when nonnegative mass vanishes.
\end{proof}
\end{corollary}

\begin{corollary} \label{cor:domination-implies-vitali-side-conditions} \hlt{Domination Gives Tail and Small-Set Control}\\
Let $1\leq p<\infty$, $g\in L^p(\mu)$, and $\abs{f_n}\leq g$ a.e. for all $n$. Then, for every $\varepsilon>0$:
\begin{enumerate}[label=(\roman*)]
\setlength{\itemsep}{0pt}
\item $\exists E_\varepsilon\in\mathcal M$ with $\mu(E_\varepsilon)<+\infty$ such that $F\cap E_\varepsilon=\emptyset\Rightarrow\sup_n\int_F\abs{f_n}^p\,d\mu<\varepsilon^p$.
\item $\exists\delta(\varepsilon)>0$ such that $\mu(E)<\delta(\varepsilon)\Rightarrow\sup_n\int_E\abs{f_n}^p\,d\mu<\varepsilon^p$.
\end{enumerate}
\begin{proof}
\begin{enumerate}[label=(\roman*)]
\setlength{\itemsep}{0pt}
\item Proposition \ref{prop:simple-density-support-and-tails-lp} gives $E_\varepsilon$ with $\mu(E_\varepsilon)<+\infty$ and $F\cap E_\varepsilon=\emptyset\Rightarrow\int_F\abs g^p\,d\mu<\varepsilon^p$.\\
Since $\abs{f_n}^p\leq\abs g^p$ a.e., $\sup_n\int_F\abs{f_n}^p\,d\mu<\varepsilon^p$.
\item Lemma \ref{lem:lp-integral-absolute-continuity} gives $\delta(\varepsilon)>0$ with $\mu(E)<\delta(\varepsilon)\Rightarrow\int_E\abs g^p\,d\mu<\varepsilon^p$.\\
Since $\abs{f_n}^p\leq\abs g^p$ a.e., $\sup_n\int_E\abs{f_n}^p\,d\mu<\varepsilon^p$.
\end{enumerate}
\hlt{Structural Commentary}: One $L^p$ envelope supplies both tail tightness and small-set control.\\
\hlt{Key Insight}: Domination transfers both estimates from $g$ to every $f_n$.
\end{proof}
\end{corollary}
